% TODO
% Response to referee 1, comment 6
Plan:
\begin{enumerate}
\item This is true.
\item We need to read these papers and understand what is different about the electromagnons the referee is referring to.
\item For example, if they're all transverse electromagnons, we can just change the descriptor in our manuscript to ``transverse''.
\end{enumerate}

We completely agree that other types of magnons can be excited via electric fields, not just pseudo-Goldstone magnons.
We have modifed the phrasing in the revised version of the manuscript.

\begin{changesmade}
Revised the paragraph as follows: ``Conventionally, the \gls{dme} is realized by coupling an external electric field to an electric-dipole-active magnon, i.e., an electromagnon.
A variety of electromagnons have been identified in multiferroic materials, including perovskite rare-earth manganites\citep{kida_terahertz_2009,tokura_multiferroics_2014}, \gls{vdw} magnets\citep{kim_terahertz_2023}, and others.
Two types of electromagnons are relevant to this work.
The first is the pseudo-\golstone magnon (Figs. 1b,c), which becomes electric-dipole active in type-II multiferroics and whose energy is set by the magnetic anisotropy gap.
The second is due to the \higgs mode of the ordered multiferroic phase (\cref{fig:fig1d,fig:fig1e}), which involves fluctuations in the \textit{magnitude} of the spins.
This mode is also, in principle, electric-dipole active in type-II multiferroics;\cite{matsumoto_electromagnon_2014} however, since it typically resides at a higher energy than the \goldstone modes, it is expected to decay quickly into the latter and remain overdamped.\citep{jain_higgs_2017}''
\end{changesmade}
